\section{RISC-V Confidential Computing}

RISC-V confidential computing is best understood as the convergence of two
historically separate lines of work. The first line is the open enclave
literature, where researchers use RISC-V's extensible ISA, privilege model, and
physical-memory protection mechanisms to build small trusted monitors and
enclave runtimes. The second line is the standardization-oriented confidential
VM effort represented by CoVE and AP-TEE, where the target is no longer only an
application enclave but a tenant VM protected from an untrusted host
environment~\cite{sahita2023cove,riscv_ap_tee_2024,boubakri2025riscvtee}.
Boubakri et al. are used here as survey/background substrate; mechanism and
status claims in this section rely on the original CoVE paper, AP-TEE draft, or
other primary sources. This distinction is important for classification:
PMP-based enclaves, tagged-memory enclaves, and CoVE-style TVMs share a
motivation, but they do not expose the same lifecycle, device, or attestation
contract.

\subsection{From Enclave Research to Confidential VMs}

\subsubsection{Early RISC-V Enclave Design}

\textbf{Sanctum} is one of the earliest and most influential examples of using
RISC-V to explore strong software isolation through minimal hardware
extensions~\cite{costan2016sanctum}. Its importance is not just that it provides
an enclave system, but that it frames open hardware as a way to make TEE
assumptions explicit. Instead of treating a proprietary platform as a fixed
black box, Sanctum exposes the processor, cache, memory, and monitor assumptions
that must hold for enclave isolation.

\textbf{Keystone} then provides a more reusable open-source RISC-V TEE
framework built around a security monitor, PMP-based memory isolation, and
runtime abstractions for enclaves~\cite{lee2020keystone}. Keystone is central to
the survey because it shows the strengths and limitations of using standard
RISC-V mechanisms for protected execution. PMP gives a simple and efficient
physical-memory filter, but a small number of PMP regions and monitor-managed
state are not the same as a cloud-scale confidential-VM memory ownership model.
This is the point at which RISC-V enclave research begins to diverge from
later CoVE/AP-TEE work.

\subsubsection{Scaling and Specializing the Enclave Boundary}

Subsequent RISC-V enclave systems explore different ways to scale, specialize,
or strengthen the basic monitor-plus-memory-protection model. \textbf{TIMBER-V}
uses tag-isolated memory to support fine-grained enclaves, which is especially
relevant for embedded and compartmentalized systems where the protected unit may
be smaller than a full VM~\cite{weiser2019timberv}. \textbf{CURE} introduces
customizable resilient enclaves, emphasizing configurable isolation and
resilience tradeoffs~\cite{bahmani2021cure}. \textbf{MI6} studies secure
enclaves in the context of a speculative out-of-order processor, making clear
that the microarchitecture matters even when the ISA-level enclave abstraction
appears simple~\cite{bourgeat2019mi6}.

\textbf{Penglai} addresses a different bottleneck: scalable memory protection
for large numbers of enclaves, motivated by cloud and serverless deployment
patterns~\cite{feng2021penglai}. \textbf{SPEAR-V} proposes low-overhead
enclave primitives for RISC-V and is useful as a later point in the lineage
because it studies practical primitive design rather than only monitor
structure~\cite{schrammel2023spearv}. \textbf{Cerberus} focuses on formal and
efficient enclave memory sharing, an issue that becomes increasingly important
once protected domains need deliberate shared-memory communication rather than
complete isolation~\cite{lee2022cerberus}. \textbf{ACE} extends the discussion
toward confidential computing for embedded RISC-V systems~\cite{ozga2025ace}.

Across these works, the recurring lesson is that a TEE is not only an execution
mode. It is a contract over memory ownership, entry and exit, metadata
integrity, sharing, interrupt behavior, and attestation. Earlier enclave
systems often solve a subset of that contract for a specific threat model.
CoVE/AP-TEE attempts to standardize a larger contract for confidential VMs.

\subsection{CoVE and AP-TEE}

\subsubsection{Threat Model and Architectural Role}

RISC-V CoVE targets a host-resilient confidential-VM model. The tenant workload
executes in a trusted virtual machine while a host stack remains responsible for
resource management, scheduling, and device orchestration~\cite{sahita2023cove}.
This is conceptually close to Arm CCA in that both designs separate management
authority from data visibility, but the RISC-V path is expressed through a
different standardization structure and a more modular ecosystem. The AP-TEE
specification defines the relevant ABI and lifecycle concepts, but it remains a
draft and should be treated as a standardization work in progress rather than a
ratified architecture~\cite{riscv_ap_tee_2024}.

The core AP-TEE objects are the \textbf{Trusted Security Monitor (TSM)}, the
\textbf{TSM driver}, \textbf{Supervisor Domains}, and the \textbf{Trusted
Virtual Machine (TVM)}~\cite{riscv_ap_tee_2024}. The host retains the ability
to allocate resources, but the TSM mediates TVM creation, memory assignment,
entry, exit, and security-sensitive transitions. This moves RISC-V beyond the
simple observation that PMP can isolate memory. A confidential VM needs a
defined lifecycle, a way to donate and reclaim pages, an attestation path, and
an interface that lets untrusted management software perform useful work without
reading protected state.

\subsubsection{Memory Lifecycle and TVM State}

Memory lifecycle is the central mechanism that differentiates CoVE/AP-TEE from
earlier PMP-only enclave explanations. The draft AP-TEE model includes memory
\textbf{donation}, \textbf{reclaim}, and \textbf{sharing} semantics, which are
needed because a TVM is dynamic: pages are created, mapped, unmapped, shared for
communication, and eventually returned to the host~\cite{riscv_ap_tee_2024}.
These transitions are security-critical. If a page can move between host and TVM
ownership without explicit state checks, stale mappings, unclean metadata, or
ambiguous sharing state can break the confidentiality boundary.

This is also where RISC-V confidential computing should not be reduced to
memory encryption. The AP-TEE draft primarily specifies ownership, lifecycle,
and access-control semantics; it does not by itself replace a memory-encryption
and replay-protection taxonomy~\cite{henson2014memory}. A complete platform may
combine access control, encryption, integrity protection, and attestation, but
those mechanisms answer different questions. Access control asks which agent
may reach a physical page. Encryption asks what an observer of raw memory can
learn. Integrity and replay protection ask whether stale or modified data can be
detected. A precise survey has to keep those layers separate.

\subsubsection{Attestation and Evidence}

AP-TEE also requires an attestation story because tenants need external
evidence before provisioning secrets to a TVM~\cite{riscv_ap_tee_2024}. The
evidence has to bind the platform, TSM, TVM initial state, and relevant
configuration to a verifier-visible claim. This mirrors the broader attestation
literature: a TEE claim is useful only if the measured state is well-defined,
fresh, tied to the execution context that will receive secrets, and appraised by
a verifier under a clear RATS-style role model~\cite{rats_rfc,eat_rfc,menetrey2022attestation}.

In RISC-V systems, this evidence chain is especially important because the
ecosystem encourages modular composition. A TVM may depend on a TSM, a platform
root of trust, an IOMMU, interrupt controllers, device-security protocols, and
firmware components from different vendors or projects. Attestation is the
mechanism that turns that composition into something a remote tenant can
evaluate.

\subsection{RISC-V Confidential I/O}

\subsubsection{Why CoVE Needs I/O Standardization}

A confidential VM that cannot use real devices is not sufficient for cloud or
edge deployment. CoVE-IO and TEE-I/O therefore extend the TVM discussion into
device identity, DMA, MMIO, interrupts, and link protection~\cite{riscv_cove_io_2026}. The
draft CoVE-IO specification is not ratified, but it is already useful as a map
of the required mechanisms. It introduces concepts such as trusted device
interfaces and device/security managers while relying on lower-level building
blocks including SPDM, TDISP, PCIe IDE, IOMMU support, and trusted interrupt
delivery~\cite{riscv_cove_io_2026,dmtf_spdm_2025,pcisig_tdisp_2022,pcie_ide}.
The TDISP citation is public PCI-SIG metadata for a member-gated ECN, so it
supports only the device-interface lifecycle role unless a stronger public
source is cited for a specific normative rule.

The design implication is straightforward: CoVE cannot stop at CPU-side page
ownership. A device assigned to a TVM must be identified, measured or certified,
placed into an appropriate trust state, constrained in the memory it can access,
and connected to the TVM through an interrupt path that does not silently return
control to the untrusted host. If any of those steps is missing, the I/O path
becomes a privileged channel around the confidential-computing boundary.

\subsubsection{RISC-V IOMMU, IOPMP, AIA, and sIOPMP}

The RISC-V IOMMU and AIA specifications provide two of the required supporting
layers. The IOMMU constrains device-side address translation and access
permissions, while AIA provides the interrupt and MSI architecture that CoVE-IO
can build on for trusted delivery~\cite{riscv_iommu_2023,riscv_aia_2023}. These
specifications are not confidential-computing systems by themselves. Their role
is to provide the architectural substrate on which a confidential I/O protocol
can express and enforce device ownership.

IOPMP and sIOPMP address a related but distinct problem: scalable filtering of
I/O transactions so that requester identity and memory-domain policy can be
checked outside the CPU's ordinary memory-access path~\cite{riscv_iopmp_2026,feng2024siopmp}.
For RISC-V confidential computing, this matters because DMA is the obvious way
to bypass a CPU-centered protection story. A TVM memory page protected from host
loads and stores is still exposed if a bus master can reach it without going
through a domain-aware check.

\subsection{Comparison with Arm CCA}

\begin{table}[htbp]
\centering
\caption{Arm CCA and RISC-V CoVE/AP-TEE Design Mapping}
\label{tab:arm_cca_riscv_cove}
\begin{tabular}{@{}p{0.27\linewidth}p{0.31\linewidth}p{0.31\linewidth}@{}}
\toprule
\textbf{Design Point} & \textbf{Arm CCA} & \textbf{RISC-V CoVE/AP-TEE} \\ \midrule
Protected workload & Realm & Trusted Virtual Machine \\
Management component & RMM & TSM and TSM driver \\
Host role & Resource management outside the Realm data boundary & Resource management outside the TVM data boundary \\
Memory ownership & GPT/GPC and Realm lifecycle & TVM memory donation, reclaim, sharing, and TSM-mediated state \\
Device path & SMMU and realm-aware device mediation & CoVE-IO, IOMMU, IOPMP, AIA, SPDM, TDISP, PCIe IDE \\
Standard status & Arm architecture/specification path & AP-TEE and CoVE-IO drafts, not ratified \\ \bottomrule
\end{tabular}
\end{table}

Table~\ref{tab:arm_cca_riscv_cove} shows why the comparison should be made at
the confidential-VM layer rather than between CCA and bare PMP. PMP remains a
useful primitive, especially for monitors and embedded TEEs, but CoVE/AP-TEE is
the closer architectural analogue to Arm CCA. Both systems try to let an
untrusted host manage resources without giving it data visibility. The
difference is that Arm presents a more vertically integrated Realm architecture,
whereas RISC-V standardization exposes a set of modular pieces that platform
builders must assemble into a coherent protection boundary.

\subsection{Summary}

RISC-V confidential computing has a rich pre-standardization lineage in
open-source enclave systems and a rapidly developing standardization path in
CoVE/AP-TEE. Earlier systems such as Sanctum, Keystone, TIMBER-V, CURE, MI6,
Penglai, SPEAR-V, Cerberus, and ACE establish the design vocabulary for
isolation, monitor structure, metadata, and memory sharing. CoVE/AP-TEE raises
the abstraction to confidential VMs, with TVM lifecycle, TSM mediation, memory
ownership transitions, and attestation as first-class concerns. CoVE-IO then
extends the same reasoning to devices and networks. The resulting lesson is
that RISC-V confidential computing is not one mechanism but a platform
composition problem: CPU privilege, memory ownership, I/O translation,
interrupt routing, device identity, and attestation all have to agree on the
same security boundary.
